\part{Composite Relaxations---\hspace{0pt}Admissible Records and the
Interaction Matrix}

\secn{44}{44}{Overview of Part VI}

Part VI studies two relaxations together: admissible registration (Part II) and anchoring (Part IV). The device that makes this possible is simple. In the exact setting, the presentation, the context, and the verdict of an anchored view are themselves views on the contrast domain, so the records of Part IV become compound views of an extended view family, and the whole of Parts II and V applies to them unchanged.

Section 45 develops the exact record calculus: the attribution processing inequality and the anatomy of leaks (Proposition 45.2), and the coherence law, by which admissible readings are preserved along every map between records exactly when the admissibility structure is monotone and insensitive to redundant views (Theorem 45.4). Section 46 applies the interval calculus to covers that mix feature views and record views. No new law appears (Proposition 46.2), but the policy becomes a parameter that creates and destroys cross-layer anomalies (Theorem 46.3), and sterile records reduce to their presentations (Proposition 46.4).

Sections 47--48 add noise. They give sufficient conditions under which corruption leaves admissible content unchanged---\hspace{0pt}fixed decoders under blind contamination (Proposition 47.1), decoders that cannot see the clutter's variation (Theorem 47.3)---\hspace{0pt}show that demanding immunity to all contamination forces a null decoder (Proposition 47.4), and show that budget brackets and loyalty extremality survive under stated hypotheses (Propositions 48.1, 48.2). Section 50 proves recovery and collects the pairwise interactions of the three relaxations in a single matrix.

\secn{45}{45}{Admissible records: the extended family and the coherence law}

\subsecn{1}{45.1}{Records decompose into views}

\begin{thmdef}[\hypertarget{stmt:45.1}{Definition~45.1}\ {\normalfont\upshape(record resolution; the extended family)}]

Let \(v\) be an exact anchored view: a deterministic emission \(e_v : D \to Y_v \times C_v\) with an \(\eta_v\)-compatible selection policy \(a_v\) (Definitions 19.5--19.6 with (A2)). Its \textbf{atomic record views} are the deterministic views
\[
N_v = \mathrm{pr}_Y \circ e_v : D \to Y_v,
\qquad
\chi_v = \mathrm{pr}_C \circ e_v : D \to C_v,
\qquad
\tau_v = a_v \circ e_v : D \to \widehat{T}
\]
(presentation, context, verdict). The records of \hyperlink{stmt:27.6}{Definition~27.6} are their compounds: \(S_v \cong \langle \{N_v, \tau_v\} \rangle\) and \(F_v \cong \langle \{N_v, \chi_v, \tau_v\} \rangle\). The \textbf{extended family} is \(\widehat{V} = V \cup \{N_v, \chi_v, \tau_v : v \text{ anchored}\}\), and all Part II and Part V structure---\hspace{0pt}admissibility structures, reading classes, covers, defects---\hspace{0pt}is taken over \(\mathcal{F}(\widehat{V})\).
\end{thmdef}

Nothing needs to be re-founded: the atomic record views satisfy (P2), the compounds are ordinary compound views, and the kernel operates on \(\widehat{V}\) as it does on \(V\). The first consequence is the exact record calculus.

\begin{thmplain}[\hypertarget{stmt:45.2}{Proposition~45.2}\ {\normalfont\upshape(exact record calculus; the leak anatomy)}]

For an exact anchored view \(v\):

\begin{enumerate}
\def\labelenumi{\arabic{enumi}.}
\tightlist
\item
  \textbf{(Exact attribution processing inequality.)} \(\ker(\tau_v) \le \ker(N_v) \vee \ker(\chi_v)\), since \(\tau_v\) factors through \(e_v\) (Lemma 14.2); hence \(\ker(F_v) = \ker(N_v) \vee \ker(\chi_v)\), and adjoining the verdict to any set containing \(N_v\) and \(\chi_v\) changes no \(\sigma\). This is the exact shadow of \hyperlink{stmt:27.8}{Proposition~27.8}: attribution is informationally free given the full raw presentation.
\item
  \textbf{(Leakage, exactly.)} \(\ker(S_v) = \ker(N_v) \vee \ker(\tau_v)\), and \(v\) \textbf{leaks} (under its policy) iff \(\ker(\tau_v) \not\le \ker(N_v)\), iff \(\ker(S_v) > \ker(N_v)\)---\hspace{0pt}the exact form of Proposition 28.2's phenomenon.
\item
  \textbf{(Anatomy.)} Call a pair \(Z, Z'\) \textbf{critical} if \(N_v(Z) = N_v(Z')\) and \(\chi_v(Z) \ne \chi_v(Z')\), and write \(e_v(Z) = (y, c)\), \(e_v(Z') = (y, c')\). Then exactly one of three cases holds:

  \begin{itemize}
  \tightlist
  \item
    \textbf{(taut)} \(\eta_v(y,c) = \eta_v(y,c')\) is a singleton: no compatible policy separates the pair through its verdict;
  \item
    \textbf{(forced)} \(\eta_v(y,c) \cap \eta_v(y,c') = \varnothing\) (in particular, distinct singletons---\hspace{0pt}e.g.~one sound, one failed status): \emph{every} compatible policy separates it;
  \item
    \textbf{(optional)} otherwise (a shared option exists, and at least one side has room): separation is policy-dependent.
  \end{itemize}

  Some policy leaks at \(v\) iff some critical pair is forced or optional. The all-policies criterion, however, is \textbf{global over presentation fibers, not pairwise}: for a realized presentation \(y\) let \(C_y = \{c : (y, c) = e_v(Z) \text{ for some } Z \in D\}\); a compatible policy is leak-free iff its verdict is constant on every realized fiber, which is possible iff
  \[
  \bigcap_{c \in C_y} \eta_v(y, c) \;\ne\; \varnothing \qquad \text{for every realized } y.
  \]
  Hence \textbf{every compatible policy leaks iff some realized fiber has empty total intersection}. A forced pair---\hspace{0pt}empty \emph{pairwise} intersection---\hspace{0pt}is sufficient but not necessary: on fibers with at most two realized contexts the pairwise and global criteria coincide, and on larger fibers they separate. With \(C_y = \{c_1, c_2, c_3\}\) and \(\eta(y, c_1) = \{t_1, t_2\}\), \(\eta(y, c_2) = \{t_2, t_3\}\), \(\eta(y, c_3) = \{t_1, t_3\}\), every critical pair is optional, yet the total intersection is empty, so every compatible policy separates some pair---\hspace{0pt}a Helly-type failure of the pairwise reading. Forced leaks are correspondence-driven---\hspace{0pt}Proposition 28.2's mechanism, where the anchoring \emph{status} varies across the pair---\hspace{0pt}optional leaks are policy-driven---\hspace{0pt}Proposition 30.2's mechanism---\hspace{0pt}and the fiberwise criterion exhibits a third possibility the pairwise anatomy cannot: a leak forced globally by a correspondence no two of whose values conflict. The erroneous/irreparable distinction of \hyperlink{stmt:19.5}{Definition~19.5} thereby refines to the leak.
\item
  \textbf{(Exact sterility.)} If \(\ker(\chi_v) \le \ker(N_v)\) (context draws no distinctions beyond the presentation---\hspace{0pt}the exact form of \hyperlink{stmt:30.3}{Definition~30.3}), there are no critical pairs, hence no policy leaks: \(\ker(\tau_a) \le \ker(N_v)\) for every compatible policy \(a\). More exactly: every policy is leak-free at \(v\) iff every critical pair is taut; sterility is the case of no critical pairs at all.
\end{enumerate}
\end{thmplain}

\begin{proof}

1 and 2 are \hyperlink{stmt:14.2}{Lemma~14.2} and \hyperlink{stmt:3.3}{Fact~3.3} applied to the stated factorizations. 3: on a critical pair, \(\tau\) separates iff \(a(y,c) \ne a(y,c')\); compatibility confines the two values to the two \(\eta\)-sets, so equality is impossible iff the sets are disjoint, avoidable iff a common element can be chosen, and necessarily equal iff both sets are the same singleton; a non-critical \(N\)-equal pair has \(e_v(Z) = e_v(Z')\) and hence equal verdicts under any policy, so all leakage lives on critical pairs, and the pairwise statements follow. For the global criterion: leak-freedom is exactly verdict-constancy on every realized fiber, a constant admissible choice on the fiber of \(y\) is exactly an element of \(\bigcap_{c \in C_y} \eta_v(y,c)\), and any family of such elements assembles into a compatible leak-free policy (off-fiber values chosen arbitrarily within \(\eta\)); the displayed three-context configuration witnesses that the pairwise condition is strictly weaker. 4: with no critical pairs, \(N\)-equal implies \(e_v\)-equal implies \(\tau\)-equal for every policy, i.e.~\(\ker(\tau_a) \le \ker(N_v)\); the refinement restates 3.
\end{proof}

\subsecn{2}{45.2}{Determination restriction and the coherence theorem}

The ladder maps \(F_v \twoheadrightarrow S_v \twoheadrightarrow N_v\) are not coordinate projections between view \emph{sets}; they are surjections \emph{within} single views. Does Proposition 14.4′'s closure notion extend to them? It does---\hspace{0pt}and the correct generality is not the ladder but determination itself.

\begin{thmdef}[\hypertarget{stmt:45.3}{Definition~45.3}\ {\normalfont\upshape(determination restriction; determination-monotone
structures)}]

Let \(T', T \subseteq \widehat{V}\) be finite with \(\sigma(T') \le \sigma(T)\). Then the compound \(c_T\) determines \(c_{T'}\): the map \(m_{T',T}\) defined on realized profiles by \(m_{T',T}(c_T(Z)) = c_{T'}(Z)\) is well defined (equal \(T\)-profiles lie in one \(\sigma(T)\)-block, hence in one \(\sigma(T')\)-block, hence have equal \(T'\)-profiles), extended arbitrarily off the image. The family of reading classes \(\{\mathcal{K}(T)\}_{T \in \mathcal{F}(\widehat V)}\) (Proposition 14.4′) is \textbf{closed under determination restriction} if \(\sigma(T') \le \sigma(T)\) and \(\kappa' \in \mathcal{K}(T')\) imply \(\kappa' \circ m_{T',T} \in \mathcal{K}(T)\). The structure \(K\) is \textbf{determination-monotone} if \(\sigma(T') \le \sigma(T) \Rightarrow \gamma_K(T') \le \gamma_K(T)\).
\end{thmdef}

For \(T' \subseteq T\) the map \(m_{T',T}\) is the coordinate projection (up to profile identification), so determination restriction extends Proposition 14.4′'s notion; and the record ladder is the case \(T' = \{N_v\}\) or \(\{N_v, \tau_v\}\) against \(T = \{N_v, \chi_v, \tau_v\}\) and its relatives, since \(\ker(N) \le \ker(S) \le \ker(F)\) makes each ladder surjection a determination restriction between the corresponding singletons of \(\widehat V\)'s compounds.

\begin{thmplain}[\hypertarget{stmt:45.4}{Theorem~45.4}\ {\normalfont\upshape(the coherence law for records)}]

For an admissibility structure \(K\) on \(\mathcal{F}(\widehat{V})\), the following are equivalent:

\begin{enumerate}
\def\labelenumi{\arabic{enumi}.}
\tightlist
\item
  \(\{\mathcal{K}(T)\}\) is closed under determination restriction;
\item
  \(K\) is determination-monotone;
\item
  \(\gamma_K\) factors through the content map as a monotone assignment: \(\gamma_K = g \circ \sigma\) for a unique monotone \(g : \mathrm{Ach}(\widehat{V}) \to \mathrm{Part}(D)\);
\item
  \(K\) is monotone (Definition 14.3) \textbf{and} determination-stable (Definition 35.3).
\end{enumerate}
\end{thmplain}

\begin{proof}

(1)\(\Leftrightarrow\)(2): verbatim the proof of Proposition 14.4′ with \(m_{T',T}\) in place of \(\mathrm{pr}_{T'}\)---\hspace{0pt}soundness of \(\kappa' \circ m_{T',T}\) on \(c_T\) holds because \((\kappa' \circ m_{T',T})(c_T(Z)) = \kappa'(c_{T'}(Z)) \ni Z\), its registered partition is \(\ker(\kappa' \circ c_{T'}) = \rho(\kappa')\), and the two directions then read exactly as there, with \hyperlink{stmt:14.4}{Lemma~14.4} supplying the ceiling-realizing \(\kappa'\) for necessity. (2)\(\Rightarrow\)(3): \(g(\sigma(T)) := \gamma_K(T)\) is well defined on \(\mathrm{Im}\,\sigma = \mathrm{Ach}(\widehat V)\) (equal \(\sigma\)'s give inequalities both ways) and monotone by definition of determination-monotonicity; uniqueness is surjectivity of \(\sigma\) onto \(\mathrm{Ach}\). (3)\(\Rightarrow\)(2) is immediate. (2)\(\Rightarrow\)(4): inclusions and \(\sigma\)-equalities are special determinations. (4)\(\Rightarrow\)(2): if \(\sigma(T') \le \sigma(T)\) then \(\sigma(T \cup T') = \sigma(T)\), so stability gives \(\gamma_K(T \cup T') = \gamma_K(T)\) and monotonicity gives \(\gamma_K(T') \le \gamma_K(T \cup T') = \gamma_K(T)\).
\end{proof}

\begin{thmplain}[\hypertarget{stmt:45.5}{Corollary~45.5}\ {\normalfont\upshape(record coherence and repackaging)}]

Reading classes are closed along \textbf{all determinations} iff \(K\) is determination-monotone, equivalently monotone and determination-stable. Invariance ceilings \(P_G\wedge\sigma(T)\) satisfy this condition. Closure along nested sets of record atoms alone follows already from monotonicity.

Let \(\mathcal N\) be the set of presentation atoms \(N_v\). The presentation-only structure \(\gamma(T)=\sigma(T\cap\mathcal N)\) is monotone. It can fail all-determination closure when a non-presentation view \(\tau\) determines a nonconstant presentation \(N'\): \(\ker N'\le\ker\tau\) but \(\gamma(\{N'\})=\ker N'>\bot=\gamma(\{\tau\})\). Thus its possible incoherence concerns repackaging, not the nested atomic ladder. If these kernels are equal, it directly violates stability; if they are strictly ordered, monotonicity plus \hyperlink{stmt:45.4}{Theorem~45.4} implies failure of stability somewhere in the family.
\end{thmplain}

\begin{proof}

The equivalence and last inference are \hyperlink{stmt:45.4}{Theorem~45.4}. Meets are monotone, giving the invariance example. Intersecting nested view sets with \(\mathcal N\) preserves their inclusion, proving the presentation-only monotonicity and the stated counterexample.
\end{proof}

\hyperlink{stmt:45.4}{Theorem~45.4} combines projection coherence with invariance under repackaging. Nested atomic record sets require only monotonicity; all determination maps require the additional stability condition. The extended family then lets the existing interval calculus address mixed feature/record covers.

\secn{46}{46}{Mixed covers: compounding and absorption}

\subsecn{1}{46.1}{Conservativity and the interaction identity}

\begin{thmdef}[\hypertarget{stmt:46.1}{Definition~46.1}\ {\normalfont\upshape(mixed sets; cross-layer defect)}]

For finite \(T \subseteq \widehat{V}\) write \(T_f = T \cap V\) (the feature part) and \(T_r = T \setminus T_f\) (the record part), assumed both nonempty; the \textbf{layer cover} is \(\mathcal{U}_{\mathrm{lay}}(T) = \{T_f, T_r\}\), and the \textbf{cross-layer defect} is
\[
X_K(T) \;=\; \delta_K(\mathcal{U}_{\mathrm{lay}}; T) \;=\; \big[\, \gamma_K(T_f) \vee \gamma_K(T_r),\ \gamma_K(T) \,\big].
\]
\end{thmdef}

\begin{thmplain}[\hypertarget{stmt:46.2}{Proposition~46.2}\ {\normalfont\upshape(conservativity; the interaction identity)}]

Let \(K\) be monotone on \(\mathcal{F}(\widehat V)\).

\begin{enumerate}
\def\labelenumi{\arabic{enumi}.}
\tightlist
\item
  Mixed covers are covers in \(\mathcal{F}(\widehat{V})\), so \hyperlink{stmt:37.1}{Definition~37.1}, \hyperlink{stmt:37.2}{Proposition~37.2}, \hyperlink{stmt:37.3}{Theorem~37.3}, and \hyperlink{stmt:37.4}{Lemma~37.4} apply \textbf{verbatim}: the composite theory adds no new law, and every mixed anomaly is an ordinary anomaly of the extended family.
\item
  Refining the layer cover by singletons, \hyperlink{stmt:37.4}{Lemma~37.4} gives the factorization
  \[\begin{gathered}\Delta_K(T) \;=\; J_K(T) \cdot X_K(T), \\ J_K(T) = \big[\, \sigma^{\mathrm{sep}}_K(T_f) \vee \sigma^{\mathrm{sep}}_K(T_r),\ \gamma_K(T_f) \vee \gamma_K(T_r) \,\big],\end{gathered}\]
  with \(J_K(T)\) the \textbf{join of the pure-layer anomalies} (degenerate whenever both \(\Delta_K(T_f)\) and \(\Delta_K(T_r)\) are; the converse fails, by absorption) and \(X_K(T)\) the genuinely cross-layer part. The two factors are independent coordinates of the total defect: \hyperlink{sec:46.2}{§46.2} exhibits a nondegenerate cross-layer factor with degenerate pure-layer factors, and a nondegenerate pure-layer factor absorbed by a degenerate composite.
\end{enumerate}
\end{thmplain}

\begin{proof}

1 is Definition 45.1---\hspace{0pt}nothing in \hyperlink{sec:37}{§37} used anything about \(V\) beyond its being a view family. 2: the singleton cover of \(T\) is the composite of \(\mathcal{U}_{\mathrm{lay}}\) with the singleton covers of \(T_f\) and \(T_r\); \hyperlink{stmt:37.4}{Lemma~37.4} computes the left endpoint as \(\sigma^{\mathrm{sep}}_K(T_f) \vee \sigma^{\mathrm{sep}}_K(T_r) = \sigma^{\mathrm{sep}}_K(T)\), and the concatenation is the interval identity \([a,b] = [a,m]\cdot[m,b]\) at \(m = \gamma_K(T_f) \vee \gamma_K(T_r)\), which lies between the endpoints by \hyperlink{stmt:14.6}{Proposition~14.6} applied layerwise and monotonicity. Degeneracy of \(J\) from degeneracy of both pure anomalies is the join of two equalities.
\end{proof}

\subsecn{2}{46.2}{Both interaction modes, on one configuration}

\begin{thmplain}[\hypertarget{stmt:46.3}{Theorem~46.3}\ {\normalfont\upshape(compounding and absorption)}]

There is a four-element contrast domain, an invariance structure, one feature view, and one exact anchored view such that:

\begin{enumerate}
\def\labelenumi{\arabic{enumi}.}
\tightlist
\item
  \textbf{(Compounding: the leak completes the share.)} Under the context-reading policy \(a\), the mixed pair \(T = \{v, \tau_a\}\) has both pure-layer anomalies degenerate and cross-layer defect \textbf{maximal}, \(X_{K_G}(T) = [\bot, P_G]\); under the loyal policy \(\ell\), the same pair has \(X_{K_G}(\{v, \tau_\ell\})\) degenerate. The cross-layer anomaly is \textbf{policy-created}: the loyalty anomaly of \hyperlink{stmt:30.2}{Proposition~30.2} and the witness anomaly of \hyperlink{sec:16.1}{§16.1} compound, the erring policy's verdict acting as the missing share of the parity scheme.
\item
  \textbf{(Absorption.)} On \(T = \{v, w, \tau_a\}\), the layer cover \(\{\{v,w\}, \{\tau_a\}\}\) has degenerate defect while its feature stage carries the maximal anomaly \(\Delta_{K_G}(\{v,w\}) = [\bot, P_G]\): the feature anomaly is absorbed at the layer cover. Symmetrically, the mixed cover \(\{\{v, \tau_a\}, \{w\}\}\) is degenerate while its mixed member internally carries \([\bot, P_G]\).
\end{enumerate}

Hence, asked whether the registration anomaly and the loyalty anomaly compound or absorb each other, the answer is: \textbf{both}---\hspace{0pt}the registration anomaly and the loyalty anomaly can compound, and either can absorb the other, and which occurs is a property of the cover and the policy, not of the layers. Moreover the compound instance is itself a witness anomaly of the extended family (\(v\) and \(\tau_a\) are relabeled parity shares; the witness nerve of the relevant orbit pair is \(\partial\Delta^1\)), confirming Proposition 46.2's conservativity: composition creates new \emph{instances} across layers, never a new \emph{kind}.
\end{thmplain}

\begin{proof}

Take §16.1's configuration: \(D = \{00, 01, 10, 11\}\), \(G\) the global flip with orbit partition \(P_G\) (the parity partition), invariance structure \(K_G\), feature view \(v =\) first bit, and (for item 2) \(w =\) second bit; recall \(\gamma_G(\{v\}) = P_G \wedge \ker(v) = \bot = \gamma_G(\{w\})\) and \(\gamma_G(\{v,w\}) = P_G\). The anchored view \(u\) is tracked-target (\(o \equiv t^*\)) with a distractor \(d\): \(Y_u = \{y_0\}\) (so \(N_u\) is constant), \(C_u = \{0,1\}\) with \(e_u(Z) = (y_0, b_2(Z))\), and \(\eta(y_0, 0) = \eta(y_0, 1) = \{t^*, d\}\)---\hspace{0pt}every status ambiguous, every critical pair optional (Proposition 45.2(3)). The loyal policy has \(\tau_\ell \equiv t^*\), \(\ker(\tau_\ell) = \bot\), \(\varepsilon_\ell \equiv 0\); the context-reading policy \(a\) (select \(t^*\) on \(c = 0\), \(d\) on \(c = 1\)) has \(\ker(\tau_a) = \ker(w)\) and errs with certainty on every candidate with \(b_2 = 1\)---\hspace{0pt}the exact skeleton of \hyperlink{stmt:30.2}{Proposition~30.2}.
1: For \(T = \{v, \tau_a\}\) the pure-layer anomalies are anomalies of singletons, hence degenerate. The defect: \(\gamma_G(\{\tau_a\}) = P_G \wedge \ker(w) = \bot\), and \(\gamma_G(T) = P_G \wedge (\ker(v) \vee \ker(w)) = P_G \wedge \top = P_G\); so \(X = [\bot \vee \bot, P_G] = [\bot, P_G]\), and the invariant question \(q_G\) (kernel \(P_G\)) is answerable under the ceiling from the pair jointly and from neither member---\hspace{0pt}a witness anomaly with shares \(v, \tau_a\), nerve \(\partial\Delta^1\) exactly as for \(\{v, w\}\) in Theorem 37.7's parity computation. Under \(\ell\): \(\gamma_G(\{v, \tau_\ell\}) = P_G \wedge (\ker(v) \vee \bot) = \bot\), so \(X = [\bot, \bot]\).
2: \(\gamma_G(\{v,w,\tau_a\}) = P_G \wedge \top = P_G\); the layer cover's left endpoint is \(\gamma_G(\{v,w\}) \vee \gamma_G(\{\tau_a\}) = P_G \vee \bot = P_G\): degenerate, with the feature stage anomaly \([\bot, P_G]\) intact one level down---\hspace{0pt}Lemma 37.4's absorption, now across layers. The mixed cover \(\{\{v,\tau_a\},\{w\}\}\) has left endpoint \(P_G \vee \bot = P_G\): degenerate, with its mixed member's internal defect maximal by item 1.
\end{proof}

The theorem's moral is worth separating from its arithmetic. In Part IV the loyalty anomaly was an \emph{informativeness} statement: an erring policy's record can Blackwell-dominate a faultless one's. Under admissibility it becomes an \emph{interaction} statement: the policy is a parameter that creates and destroys cross-layer obstruction classes, because the verdict view's kernel is the one datum in \(\widehat{V}\) that the interpreter's own choices set. Anchoring is thereby the layer at which the obstruction theory acquires a \emph{control} parameter---\hspace{0pt}Part V computed classes of given configurations; Part VI's configurations contain a choice, and \hyperlink{stmt:46.3}{Theorem~46.3}(1) is the statement that the choice reaches the classes.

\begin{thmplain}[\hypertarget{stmt:46.4}{Proposition~46.4}\ {\normalfont\upshape(sterile records reduce to their presentation shadow)}]

Call \(T\) \textbf{\(N\)-complete} when it contains the presentation of each anchored record atom it contains. Suppose every relevant anchored view is exactly sterile and \(K\) is determination-monotone. Delete context and verdict atoms to obtain \(\bar T\). Then \(\sigma(T)=\sigma(\bar T)\) and \(\gamma_K(T)=\gamma_K(\bar T)\). For a cover all of whose members are \(N\)-complete, its defect equals that of its presentation-shadow cover. A zero cross-layer interval additionally requires the shadow cover to glue.
\end{thmplain}

\begin{proof}

Sterility gives \(\ker\chi_v\le\ker N_v\) and \(\ker\tau_v\le\ker N_v\). Each deleted atom is therefore redundant in an \(N\)-complete set. Determination stability gives ceiling equality. Apply this to the union and to each cover member to preserve both interval endpoints. Their equality to each other is precisely shadow gluing, not a consequence of sterility alone.
\end{proof}

\hyperlink{stmt:46.4}{Proposition~46.4} gives a presentation-shadow reduction for N-complete sets and memberwise N-complete covers. It does not assert that arbitrary layer-grouping defects vanish. Compounding and absorption remain possible wherever the corresponding shadow or nonsterile configuration has a defect.

\secn{47}{47}{The graded corner I: fixed-decoder robustness}

The remaining sections compose all three relaxations, in the half where composition is tame. Throughout, admissibility constrains readings of the \textbf{perceived} data (a modelling choice: admissibility constrains what the interpreter can do with the data it actually receives), and patterns are as in \hyperlink{stmt:29.1}{Definition~29.1}.

\begin{thmplain}[\hypertarget{stmt:47.1}{Proposition~47.1}\ {\normalfont\upshape(fixed-decoder statistical invariance)}]

For candidate-blind contamination with common rate \(\varepsilon<1\) and a fixed decoder \(g\),
\[g\widetilde v=(1-\varepsilon)gv+\varepsilon gQ,\]
and \(\sigma^=(g\widetilde v)=\sigma^=(gv)\). This holds simultaneously for a fixed collection of decoders. It does not assert invariance of a decoder class defined by \(g\widetilde v\preceq_B\Gamma\), since that condition can change with the channel.
\end{thmplain}

\begin{proof}

Linearity gives the identity, and the affine map has positive coefficient on \(gv\), so two candidate rows are equal after processing iff they were equal before contamination. For the boundary, take a clean identity view, contaminate it to a nontrivial BSC, and use that BSC as budget. Identity decoding fits the corrupted budget but not the clean one.
\end{proof}

\begin{thmdef}[\hypertarget{stmt:47.2}{Definition~47.2}\ {\normalfont\upshape(equalizing readings)}]

Let a pattern have common rate \(\varepsilon(Z) \equiv \varepsilon \in (0,1)\) and candidate-dependent clutter \((Q_Z)\). A reading \(g\) \textbf{equalizes} the clutter if \(g \circ Q_Z = g \circ Q_{Z'}\) for all \(Z, Z'\)---\hspace{0pt}the clutter's candidate-dependence lies in distinctions \(g\) does not draw.
\end{thmdef}

\begin{thmplain}[\hypertarget{stmt:47.3}{Theorem~47.3}\ {\normalfont\upshape(equalization is sufficient for blind-contamination
behavior)}]

Fix a family of decoders and candidate-dependent clutter \(Q_Z\) at a common rate \(0<\varepsilon<1\). If each decoder \(g\) equalizes the clutter, put \(r_g=gQ_Z\), independent of \(Z\). Then
\[g\widetilde v_Z=(1-\varepsilon)gv_Z+\varepsilon r_g.\]
Each decoded experiment is therefore a blind garbling of its clean counterpart, preserves its equal-law partition, and has both deficiencies to that counterpart at most \(\varepsilon\). For fixed \(g,Q\) it is Blackwell-decreasing with the rate. These are sufficient robustness conclusions; they do not characterize all robust models or guarantee invariance of a moving output-budget class.
\end{thmplain}

\begin{proof}

The identity is linearity. Replace the decoded output by an independent draw from \(r_g\) with probability \(\varepsilon\) to obtain the garbling. Injectivity of the affine row map gives equality-partition preservation; coupling on the uncorrupted event gives TV and both deficiency bounds. Additional blind replacement produces every higher rate below one.
\end{proof}

\begin{thmplain}[\hypertarget{stmt:47.4}{Proposition~47.4}\ {\normalfont\upshape(distributional immunity is null)}]

Fix a nonempty clutter class \(\mathcal{Q}\) and suppose \(g \circ \tilde v = g \circ v\) for \textbf{all} patterns with clutter drawn from \(\mathcal{Q}\) and all rate profiles \(\varepsilon : D \to [0,1]\). Then \(g \circ Q = g \circ v(\cdot \mid Z)\) for every \(Q \in \mathcal{Q}\) and every \(Z\); hence \(g \circ v\) is constant in \(Z\) and \(g\) carries no content at any stratum.
\end{thmplain}

\begin{proof}

Take the rate spike \(\varepsilon = \mathbb{1}_{\{Z\}}\): then \(g \circ \tilde v(\cdot \mid Z) = g \circ Q\), and immunity forces \(g \circ Q = g \circ v(\cdot \mid Z)\); ranging over \(Z\) with \(Q\) fixed makes \(g \circ v\) constant.
\end{proof}

\hyperlink{stmt:47.3}{Theorem~47.3} supplies a sufficient fixed-decoder robustness mechanism by equalizing clutter. \hyperlink{stmt:47.4}{Proposition~47.4} addresses the stronger demand of distributional immunity against every rate profile and every clutter in the declared nonempty class; that demand forces a null decoded experiment. These are distinct quantified statements, not a general equivalence between robustness and invariance.

\secn{48}{48}{The graded corner II: accumulation and extremality under ceilings}

\begin{thmplain}[\hypertarget{stmt:48.1}{Proposition~48.1}\ {\normalfont\upshape(product preprocessing preserves the budget bracket)}]

Let \(\widehat{K}\) satisfy (K0̂)--(K3̂) on a coherent family, and corrupt each view \(v \in S\) independently by a candidate-blind pattern \((\varepsilon_v, Q_v)\)---\hspace{0pt}a jointly candidate-blind product pattern (Definition 29.7). Judge admissibility of readings of the perceived family by the \textbf{same} ceilings \(\Gamma_{\widehat K}\). Then every separable admissible reading of the perceived family is a garbling of \(\Gamma_{\widehat K}(S)\): the upper bracket holds for the unchanged budget experiments. To call these budgets achievable ceilings of the perceived family, one must additionally verify (K1̂), \(\Gamma(T)\preceq_B\widetilde P_T\).
\end{thmplain}

\begin{proof}

Write \(M_v\) for the contamination garbling of \hyperlink{stmt:29.3}{Theorem~29.3}(1), so \(\widetilde{P}_S = (\bigotimes_v M_v) \circ P_S\) by joint blindness of the product pattern. A separable admissible reading of the perceived family is \((\bigotimes_v G_v) \circ \widetilde{P}_S = (\bigotimes_v (G_v \circ M_v)) \circ P_S\)---\hspace{0pt}a separable reading of the \emph{honest} family with per-view processings \(G_v \circ M_v\). Its per-view admissibility holds because \((G_v \circ M_v) \circ v = G_v \circ \tilde v\) is admissible by hypothesis; (K3̂) on the honest family then bounds the whole reading by \(\Gamma_{\widehat K}(S)\). \(\square\)

Product preprocessing is sufficient, not necessary. Appendix D.6 extends this argument to candidate-independent mixtures whose branchwise composite decoders are admissible. Joint blindness alone does not supply those branchwise bounds.
\end{proof}

\begin{thmplain}[\hypertarget{stmt:48.2}{Proposition~48.2}\ {\normalfont\upshape(extremality survives principal ceilings)}]

In the setting of \hyperlink{stmt:30.4}{Theorem~30.4} (tracked target, sterile context, clean scene), let admissibility on record readings be of ceiling type: a reading \(g\) of a record \(R\) is admissible iff its output experiment satisfies \(g \circ R \preceq_B \Gamma\) for a fixed ceiling experiment \(\Gamma\) (the (K2̂) form). Then for every compatible policy \(A\), the admissible achievable set of \(A\)'s accepted record is contained in that of the loyal record:
\[
\{\, g \circ \tilde v_A \;:\; g \text{ admissible} \,\} \;\subseteq\; \{\, h \circ \tilde v_\ell \;:\; h \text{ admissible} \,\},
\]
so loyalty remains extremal under \textbf{every} principal graded ceiling. Inversion of the comparison requires admissible classes that are not determined by their output experiments---\hspace{0pt}for example classes that fix the format of the reading as well as its output experiment, the graded analogue of §16.2's inversion. Admissibility criteria given by a lower set of experiments (Remark 24.6) are output-determined, so the proposition covers them as well.
\end{thmplain}

\begin{proof}

\hyperlink{stmt:30.4}{Theorem~30.4} gives \(\tilde v_A = G_A \circ \tilde v_\ell\); for admissible \(g\) put \(h = g \circ G_A\), whose output experiment \(h \circ \tilde v_\ell = g \circ \tilde v_A\) is the same, hence admissible under any output-determined criterion.
\end{proof}

The graded comparisons use different hypotheses: fixed-decoder invariance in \hyperlink{stmt:47.1}{Proposition~47.1}, product or branchwise-admissible preprocessing in \hyperlink{stmt:48.1}{Proposition~48.1} and Appendix D.6, and sterile-context output-determined admissibility in \hyperlink{stmt:48.2}{Proposition~48.2}. \hyperlink{stmt:41.2}{Theorem~41.2} concerns equal-law partitions under jointly blind replacement. None of these statements implies that arbitrary corruption preserves every ceiling or every record comparison.

\secn{49}{49}{Worked example H: calibration, linkage, and the completed share}

The configuration of \hyperlink{stmt:46.3}{Theorem~46.3}, with measurement semantics. Two binary assays are run on each sample; a calibration ambiguity---\hspace{0pt}the global flip \(G\)---\hspace{0pt}makes only their \emph{agreement} meaningful, so the meaningfulness ceiling is the invariance structure \(K_G\), and the meaningful question is \(q_G\) with kernel \(P_G\) (§16.1's measurement reading, Part II).

Assay A is registered as the feature view \(v\). Assay B's value is never registered as a feature: it survives only as the \emph{context} of a linkage step---\hspace{0pt}the record is attributed to the tracked batch \(t^*\) or to a distractor batch \(d\), with both attributions admissible in every context (all statuses ambiguous; every critical pair optional, \hyperlink{stmt:45.2}{Proposition~45.2}(3)).

The loyal linkage attributes everything to \(t^*\): zero anchoring error, constant verdict, and the meaningful question is unanswerable from anything in the file---\hspace{0pt}\(\gamma_G(\{v, \tau_\ell\}) = \bot\). The context-reading linkage errs on half the candidates with certainty, and its verdict is worth a share: \(\gamma_G(\{v, \tau_a\}) = P_G\), with cross-layer defect \([\bot, P_G]\)---\hspace{0pt}the agreement of the assays, the only meaningful content in the problem, is recoverable exactly from \emph{assay A plus the pattern of linkage decisions}, and from nothing less (Theorem 46.3(1)). An auditor who reads only features and discards verdicts operates the presentation-only structure of \hyperlink{stmt:45.5}{Corollary~45.5}: the leak is closed by fiat at the cost of the entire interval, and the auditor's reading discipline is provably incoherent across the record ladder (not determination-stable). Adding assay B as a feature view exhibits absorption: over \(\{v, w, \tau_a\}\) the layer cover glues, the verdict is redundant, and the anomaly of \(\{v,w\}\) survives untouched one level down (Theorem 46.3(2)). In the graded regime, blind transcription noise on the assays changes none of the admissible statistical content (Proposition 47.1); noise whose direction depends on the batch is neutralized exactly if the calibration-invariant readings cannot see its variation (Theorem 47.3); and a laboratory that demands readings \emph{distributionally} immune to arbitrary contamination has, by \hyperlink{stmt:47.4}{Proposition~47.4}, demanded readings that measure nothing.

\secn{50}{50}{Recovery, the interaction matrix, and the discharged cell}

\subsecn{1}{50.1}{Recovery}

\begin{thmplain}[\hypertarget{stmt:50.1}{Theorem~50.1}\ {\normalfont\upshape(recovery with explicit record scope)}]

Under sound unique anchoring the verdict is forced and anchoring error is zero on realized inputs. Always \(F\simeq_B E\) and \(N\preceq_B S\preceq_B F\). In the tracked-target reading \(S\simeq_B N\); equality with the evidence record additionally follows from sterile context. Under deterministic full admissibility, every registration and mixed-cover interval is degenerate, whether or not the context is sterile.
\end{thmplain}

\begin{proof}

The record statements are \hyperlink{stmt:32.1}{Theorem~32.1}. Full admissibility gives \(\gamma(T)=\sigma(T)\) and joins of view kernels preserve unions, so every cover glues, including mixed covers. This interval collapse does not imply that the presentation determines all other records.
\end{proof}

\subsecn{2}{50.2}{The interaction matrix}

The three relaxations now have all their pairwise interactions on record; the matrix collects the verdicts, with the diagonal the pure theories.

{\def\LTcaptype{none} % do not increment counter
\begin{longtable}[]{@{}
  >{\raggedright\arraybackslash}p{(\linewidth - 6\tabcolsep) * \real{0.2105}}
  >{\raggedright\arraybackslash}p{(\linewidth - 6\tabcolsep) * \real{0.2105}}
  >{\raggedright\arraybackslash}p{(\linewidth - 6\tabcolsep) * \real{0.2895}}
  >{\raggedright\arraybackslash}p{(\linewidth - 6\tabcolsep) * \real{0.2895}}@{}}
\toprule\noalign{}
\begin{minipage}[b]{\linewidth}\raggedright
\end{minipage} & \begin{minipage}[b]{\linewidth}\raggedright
registration (A3)
\end{minipage} & \begin{minipage}[b]{\linewidth}\raggedright
determinism (A2)
\end{minipage} & \begin{minipage}[b]{\linewidth}\raggedright
anchoring (A4′)
\end{minipage} \\
\midrule\noalign{}
\endhead
\bottomrule\noalign{}
\endlastfoot
\textbf{registration (A3)} & Part II: \(\Delta_K\); separability \(=\) gluing; Part V \hyperlink{sec:37}{§37}: universal defect, vanishing theorem, orbit formula, reflection \(\cdot\) witness & \hyperlink{sec:23}{§23}: accumulation refutes the bracket; (K3̂) restores it at the price of budget semantics; budget classes are representable only as non-principal lower sets (Rem. 24.6) & \textbf{Part VI}: conservative calculus (Prop. 46.2); coherence \(\iff\) determination-monotone \(=\) monotone \(+\) stable (Thm. 45.4); compounding \emph{and} absorption, policy-created classes (Thm. 46.3); memberwise N-complete sterile-shadow reduction (Prop. 46.4) \\
\textbf{determinism (A2)} & ---\hspace{0pt} & Part III: strata; coupling fiber; Part V \hyperlink{sec:39}{§39}: Vorob'ev descent, contextual evidence / underdetermination & Part IV: fidelity trichotomy; Part V §§40--41: blind \(=\) neutral endomorphism, \(2\varepsilon\)-stable enriched cells; marginal blindness creates; targeted severs \\
\textbf{anchoring (A4′)} & ---\hspace{0pt} & ---\hspace{0pt} & Part IV: leak anatomy now exact (Prop. 45.2(3)): forced (correspondence) vs.~optional (policy) \\
\end{longtable}
}

Part VIII supplies explicit joint- and marginal-blind transfer failures and computes the policy-bit examples under their distinct corruption models. Least-transfer existence is treated as an order problem, not inferred from an antichain. Appendix D.11 gives finite certificates for adequacy of a fixed transfer against a polyhedral attack class.

Reading the matrix against the tracked coordinates: codomain and compositionality (§24), representability (Remark 24.6), fidelity (§32.2), and descent (Part V \hyperlink{sec:42.2}{§42.2}) all remain in force over \(\widehat{V}\) without modification, and Part VI adds no sixth coordinate---\hspace{0pt}its finding is that the composite theory needs none. The interactions are governed by the coordinates already installed, plus one new \emph{lever}: the policy, the single datum of the architecture set by the interpreter, whose reach into the obstruction classes is \hyperlink{stmt:46.3}{Theorem~46.3}(1).

\subsecn{3}{50.3}{Summary of the composite results}

The exact record-reading result is all-determination closure iff determination monotonicity. The graded results are fixed-decoder statistical invariance and sufficient budget-bracketing constructions; they are not an extension of that exact equivalence to every graded record class. Mixed exact defects are computed by the same interval calculus, with sterility reducing memberwise N-complete covers to their shadows.

\secn{51}{51}{Summary of Part VI}

The exact part proves that records are ordinary views of an extended family, characterizes coherence of admissible readings along all record maps (Theorem 45.4), and exhibits compounding and absorption of the registration and loyalty anomalies on one configuration (Theorem 46.3), with the policy as the interpreter's control over cross-layer anomalies.

The graded results are sufficient conditions, not characterizations: fixed-decoder invariance and clutter equalization give robustness under blind or equalized contamination (Proposition 47.1, \hyperlink{stmt:47.3}{Theorem~47.3}); product preprocessing preserves the budget bracket (Proposition 48.1), with an extension to branchwise-admissible mixtures in Appendix D.6; and loyalty extremality survives admissibility criteria determined by output experiments (Proposition 48.2). None of these makes the honest budgets achievable from corrupted data, and Part VIII shows how they fail without their hypotheses.
